\documentclass{article}
\usepackage{amsmath,amssymb,hyperref}
\title{Changing interaction shapes in native synchronization}
\date{15 November 2026}
\begin{document}\maketitle
\paragraph{Status and relation to prior work.} General compatibility result, independently reviewed in the stated scope. The project's Diagram\_Lifting\_of\_Presented\_Native\_Interaction.tex already proves restriction of its presented interaction along every shape functor, retaining full native cells. We do not rediscover that operation or claim a new Kan-extension theorem. The additional question here is how the constructed synchronization evaluation and its directed assembly comparisons behave under restriction, including natural transformations between shape maps.

\paragraph{Data.} Let $\mathcal D$ be a native hom bicategory with the selected incoming lax bilimits for the small ordinary shapes under consideration, and write $\mathcal D^J=\operatorname{Ps}(J,\mathcal D)$ and $L_J:\mathcal D^J\to\mathcal D$ for the chosen synchronization pseudofunctor. A chosen cone on $F:J\to\mathcal D$ has legs $p_j:L_JF\to Fj$ and cells $F(a)p_j\Rightarrow p_k$ for $a:j\to k$. These cells need not be invertible. Its representing cone functor is an equivalence of categories, pseudonatural in the vertex. Fix all the chosen representing comparisons; do not assert literal uniqueness of cone representatives.

\paragraph{Shape arrows.} For $u:I\to J$, restrict the $F$ cone to its $u(i)$ legs. Universality gives a native arrow
\[
 w_{u,F}:L_JF\longrightarrow L_I(Fu)
\]
and an invertible modification identifying its projected cone with the restricted cone. Naturality on pseudonatural transformations $F\Rightarrow G$ and modifications between them follows by the same cone construction, yielding a pseudonatural $w_u:L_J\Rightarrow L_Iu^*$. No cofinality, adjoints of $u$, or limit preservation by another operation is assumed. There is no inferred reverse arrow.

For $v:H\to I$, the two arrows $w_{uv,F}$ and $w_{v,Fu}w_{u,F}$ project to the same restricted cone on $Fuv$. Full faithfulness gives their invertible native two-cell comparison; identity shape maps are treated similarly. Threefold composition and unit equations follow by projection. These statements include the actual pseudofunctor and restriction frames rather than equating arbitrary choices of $w$.

\paragraph{Shape two-cells are directed.} Let $\eta:u\Rightarrow v$ be a natural transformation of functors $I\to J$. Postcomposition gives the pseudonatural transformation $F\eta:Fu\Rightarrow Fv$. There is a canonical native two-cell
\[
 \omega_{\eta,F}:L_I(F\eta)\,w_{u,F}\Longrightarrow w_{v,F}.
\]
At $i\in I$ its projection is the original lax-cone cell
$F(\eta_i)p_{u(i)}\Rightarrow p_{v(i)}$, with the selected comparison frames. Naturality of $\eta$ and the cone equations prove that these components form a modification of cones; fullness gives $\omega$. It may be \emph{noninvertible}. This is the important distinction from the invertible comparisons for composition of shape functors. Identity and vertical composition of $\eta$ follow from the lax-cone unit/composition equations. Whiskering and horizontal composition of shape transformations are checked by projecting to the corresponding composites of their component arrows; the pseudofunctor constraints of $F$ and the same cone equations give their compatibility with the shape-arrow composition comparisons. Naturality in $F$ uses the full native diagram transformation and its modification equations. No truncation of a shape two-cell to an equality is used.

\paragraph{Compatibility with finite native operations.} Let $p$ be an admitted finite-arity native operation and $d_p^J:p(L_JF_i)\to L_Jp(F_i)$ its constructed synchronization comparison. Pointwise restriction commutes with $p$ through its specified pseudofunctor frames. The two routes
\[
 p(L_JF_i)\xrightarrow{d_p^J}L_Jp(F_i)
 \xrightarrow{w_u}L_I(p(F_i)u)
\]
and
\[
 p(L_JF_i)\xrightarrow{p(w_{u,F_i})}p(L_I(F_i u))
 \xrightarrow{d_p^I}L_Ip(F_i u)
\]
have an invertible native two-cell comparison. After projection to $i\in I$ both are the pointwise operation on the same restricted cone legs; the required identifications use the existing $d$ and $w$ cone frames. Full faithfulness constructs the comparison and proves its operation-substitution and shape-composition equations. For $\eta:u\Rightarrow v$, the compatibility equation involving the generally noninvertible $\omega_\eta$ follows by projection from $p$ applied to the tuple of original cone cells and the pseudofunctoriality of $p$. It is an equation of the appropriate two-cell pastes, not a claim that $\omega_\eta$ inverts.

The same argument handles admitted transformations and modifications between operations, including a directed interchanger only when its full native admission has been verified. Hence varying the interaction architecture does not require a new global assembly law beyond the bounded native operation data already used in synchronization.

\paragraph{Canonical fixed-state feedback.} In the completed pure/process doctrine, canonical flattened $B_X$ for a fixed pure state preserves the selected cones by the separate pure assembly calculation. Its comparison $b_X:B_XL_J\simeq L_JB_X$ is therefore invertible. The preceding unary-operation compatibility identifies its two pastes with $w_u$, and gives the equation with each directed $\omega_\eta$. Canonical joining and units respect these comparisons by the same projected native naturality. This applies to the reviewed fixed-$X$ operation only; no preservation of a varying shared-state tensor is inferred.

\paragraph{Scope.} The shape specifies an admitted native diagram, not arbitrary data lacking its required body squares. Shape restriction forgets or reorganizes that specified diagram; it does not reconstruct missing dependencies or furnish a physical implementation of rewiring. External process interfaces remain fixed in each hom diagram. No Kan pushforward, arbitrary shape-changing content optransport, instantaneous feedback, or reverse synchronization comparison is asserted. The mechanism is established cone restriction and its universal-property coherence, integrated here with the native synchronization and fixed-state feedback constructions. The exact-review target is the directed $\omega$ inventory and its compatibility with operation and shape compositions, not a new priority claim for restriction itself.
\end{document}
